\chapter{Síntesis científica y perspectivas}
\label{ch:conclusions}

\section{Resultados establecidos}

La evidencia disponible permite separar tres niveles que antes aparecían mezclados. Primero, la calibración B0--B2 mejora propiedades macroscópicas del canal: potencia, magnitud y dispersión de retardos. Segundo, la discrepancia compleja contiene una pendiente espectral de baja dimensión que puede predecirse parcialmente desde posición y descriptores de trazado. Tercero, las representaciones relativas eliminan grados de libertad no disponibles y mejoran la discriminabilidad espacial estática.

S4.7-R confirma este tercer nivel dentro de dc01: $U(1)$ por arreglo aumenta su victoria global de aproximadamente $0.591$ a $0.796$ después de la corrección con retardo predicho, y la fusión congelada con el producto espectral alcanza aproximadamente $0.807$. La mejora se concentra en separación regional; no aparece una ganancia local concluyente. El umbral sostenido de $1$--$2\lambda$ describe discriminabilidad del protocolo y no resolución física universal.

\section{Resultados negativos y límites identificados}

La fase común al arreglo se estima a posteriori, pero S4.7-Ua no encuentra una mejora robusta al modelar $p(\psi\mid x_{\rm RT})$ frente a distribuciones circulares incondicionales. Retirar $\beta$, restringir casos con $\kappa>0$ o introducir una escala oráculo RT no cambia materialmente la conclusión. El mínimo $\Delta\mathrm{NLL}$ negativo observado, del orden de $-2\times10^{-6}$, es despreciable.

Por tanto, bajo dc01/B0, las covariables y particiones evaluadas, $\psi$ se trata como componente de $Z_\phi^{\rm nuisance}$. Esta decisión no afirma que la fase común sea universalmente aleatoria o impredecible. Una referencia estable o metadatos distintos podrían modificar su identificabilidad.

La auditoría S4.7-Ub.0 valida la suma de trayectorias con error aproximado de $2.93\times10^{-7}$, pero revela condición muy distinta entre transmisores: Tx0 tiene rango $189/189$, $\kappa\approx2.79\times10^2$ y coherencia máxima $0.997073$; Tx1, rango $178/178$, $\kappa\approx1.45\times10^6$ y coherencia $0.999974$. Las ecuaciones normales agravan Tx1 y un ridge exploratorio introduce sesgo excesivo. El rango algebraico completo no implica identificabilidad robusta de coeficientes pathwise.

\section{Hipótesis revisadas por la evidencia}

La hipótesis de una fase absoluta directamente predecible desde geometría se ha debilitado. La evidencia favorece una estructura operacional: componente espectral parcialmente transferible, fase común tratada como perturbación bajo el protocolo actual y estructura relativa aprovechable. La factorización latente se conserva como guía conceptual, sin prometer identificación separada de todos sus bloques.

La distinción entre fidelidad estática y diferencial adquiere prioridad. S4.7-R demuestra información espacial estática, pero no respuesta a movimiento. La evidencia d010 confirma que la sensibilidad depende de geometría e infraestructura distribuida; sin embargo,
\begin{equation}
\frac{\partial H}{\partial p_{\rm Tx}}
\ne
\frac{\partial H}{\partial q_{\rm humano}},
\end{equation}
por lo que no sustituye una intervención humana o mecánica con terminales fijos.

\section{Experimentos con mayor valor informativo}

La siguiente evidencia decisiva procede de E0 y E1. E0 medirá el piso instrumental, reinicios y referencia. E1 impondrá desplazamientos conocidos y comparará canal complejo, magnitud, representación relativa, descriptor espectral y rama coherente mediante $\Delta H$, $J_{\rm real}$ y $J_{\rm sim}$. E2 utilizará un fantoma sólo después de cerrar ese régimen.

En paralelo, S4.7-Ub debe estudiar espectros singulares, rango efectivo, coherencia, separaciones de retardo y agrupación de modos no resolubles. Con $B\approx\SI{50}{MHz}$, $1/B\approx\SI{20}{ns}$ explica por qué varias firmas pueden ser casi indistinguibles. No se presentará un posterior independiente por trayectoria hasta definir un estimando resoluble. La comparación NeRF$^2$/gemelos aprendidos y la auditoría OpenLoRa avanzan en paralelo sin bloquear la física controlada.

\section{Agenda doctoral y proyecto}

La tesis busca establecer qué representación preserva sensibilidad, cuándo un gemelo reproduce esa sensibilidad y cómo la observabilidad apoya transferencia. El proyecto financiado aporta plataforma, metrología, referencias, extensión multirradio y sitio piloto. Su transición sigue
\[
\begin{aligned}
\text{instrumentación estable}&\rightarrow
\text{movimiento controlado detectable}\\
&\rightarrow\text{respiración controlada recuperable}
\rightarrow\text{persona}\rightarrow\text{piloto}.
\end{aligned}
\]
El manuscrito permanece abierto en los puntos que requieren medición. Su contribución actual es delimitar con precisión qué parte del canal estático se ha modelado, qué grados de libertad no mostraron predictibilidad y por qué la fidelidad diferencial debe gobernar la siguiente etapa.
